\documentclass[../main.tex]{subfiles}

\begin{document}

\section{Derivacao e Integracao Numerica}

\subsection{Derivadas numericas}

Foram implementadas aproximacoes progressiva, regressiva e central para a
primeira derivada. Para a segunda derivada foi usada a formula central:
\[
    f''(x) \approx \frac{f(x+h)-2f(x)+f(x-h)}{h^2}.
\]

Tambem foi incluida a extrapolacao de Richardson, que combina duas aproximacoes
com passos \(h\) e \(h/2\):
\[
    R = A(h/2) + \frac{A(h/2)-A(h)}{2^p-1},
\]
onde \(p\) e a ordem do erro do metodo usado.

\subsection{Regras de integracao}

As regras implementadas foram:

\begin{itemize}
    \item trapezio simples e composto;
    \item Simpson \(1/3\) simples e composto;
    \item Simpson \(3/8\) simples e composto;
    \item quadratura de Gauss-Legendre com 2, 3 e 4 pontos.
\end{itemize}

Para os exercicios que pedem erro relativo menor que uma tolerancia, foi usado
um processo simples: calcula-se Simpson \(1/3\) composto, dobra-se o numero de
subintervalos e compara-se a variacao relativa entre duas aproximacoes
consecutivas.

\end{document}
